\section{The exact coefficient law and the resonant scale}
\label{sec:model}

\subsection{Quadratic feedback as a compound-Poisson coefficient}

For real $s>2$ and $b>0$, set
\begin{equation}\label{eq:parameters}
 u=\sqrt{b^2+2\zeta(s)},\qquad
 A=\zeta(s-1),\qquad v_c=u-b,\qquad R_b=e^{-v_c}.
\end{equation}
Consider the unique germ
\begin{equation}\label{eq:model}
 uv(r)-\frac{v(r)^2}{2}=\Li_s(re^{v(r)}),\qquad v(0)=0,
 \qquad a_n=[r^n]v(r).
\end{equation}
Define
\begin{equation}\label{eq:f}
 f(z)=u-\sqrt{u^2-2\Li_s(z)}=\sum_{k\ge1}c_kz^k,
\end{equation}
using the square root equal to $u$ at zero. Since
$2\zeta(s)<u^2$, the Taylor expansion of the outer square root converges
absolutely on the whole image of $|z|\le1$. In particular,
\begin{equation}\label{eq:positive-intensity}
 c_k\ge\frac1{u k^s}>0,\qquad
 \sum_{k\ge1}c_k=f(1)=v_c,\qquad f'(1)=\frac Ab.
\end{equation}

\begin{proposition}[Exact coefficient--probability identity]
\label{prop:poisson}
Let $X_n$ be the nonnegative integer-valued compound-Poisson variable
with probability generating function
\begin{equation}\label{eq:pgf}
 \E z^{X_n}=\exp\{n(f(z)-f(1))\}.
\end{equation}
Then
\begin{equation}\label{eq:probability-identity}
 R_b^n a_n=\frac1n\Prob(X_n=n),\qquad
 \E X_n=\frac{nA}{b}.
\end{equation}
The identity holds for every $b>0$, including $b<A$.
\end{proposition}

\begin{proof}
Put $z=re^{v}$ in \eqref{eq:model}. Solving the quadratic equation
on its germ gives $v=f(z)$, so $z=re^{f(z)}$. Lagrange inversion gives
\[
 a_n=\frac1n[z^{n-1}]f'(z)e^{nf(z)}
     =\frac1n[z^n]e^{nf(z)}.
\]
For the second equality differentiate $e^{nf(z)}$ and compare its
$z^{n-1}$ coefficient. Positivity and summability in
\eqref{eq:positive-intensity} make
$\exp\{n\sum c_k(z^k-1)\}$ a probability generating function:
take a Poisson number of jumps of total intensity $nf(1)$, with
jump probabilities $c_k/f(1)$. Multiplication by $e^{-nf(1)}$
proves the coefficient identity. Differentiation at $z=1$ proves
the mean formula.
\end{proof}

The endpoint reference $R_b$ is the actual Taylor radius at and above
$b=A$. Below that threshold the branch first encounters an interior
fold and the radius is larger than $R_b$ \cite{repo-boundary}.
No radius claim is needed for \eqref{eq:probability-identity}.

\subsection{Pairing the colliding Gamma and zeta terms}

Write $s=3+\eps$ and define
\begin{align}
 E_2(\eps)&=\frac{\Gamma(1-\eps)}{(1+\eps)(1+\eps/2)},
 & Z(\eps)&=\zeta(1+\eps)-\eps^{-1},\label{eq:E-Z}\\
 K_\eps(\ell)&=Z(\eps)+
      \frac{1-E_2(\eps)e^{-\eps\ell}}{\eps}.
      \label{eq:K}
\end{align}
Both expressions are interpreted by removal of the singularity at
$\eps=0$. The elementary values are
\begin{equation}\label{eq:resonant-values}
 E_2(0)=1,\quad E_2'(0)=\gamma-\frac32,\quad
 Z(0)=\gamma,\quad K_0(\ell)=\ell+\frac32.
\end{equation}
Gamma recurrence gives
$2\Gamma(-2-\eps)=-E_2(\eps)/\eps$. Pairing that term with the
zeta pole in the convergent Jonqui\`ere expansion yields the following
uniform local statement.

\begin{lemma}[The centered local exponent]
\label{lem:local}
There are $r>0$ and $C<\infty$, uniform for $b\in[b_-,b_+]$ and
real $|\eps|\le1/4$, such that, for
$0<|t|<r$ and $|\arg t|\le\pi/2$,
\begin{equation}\label{eq:local-exponent}
 f(e^{-t})-f(1)+\frac Ab t
 =\frac{t^2}{2b}\left\{K_\eps(-\Log t)+\nu\right\}
   +\mathcal E(t),\qquad \nu=\frac{A^2}{b^2},
\end{equation}
where $|\mathcal E(t)|\le C|t|^{5/2}$.
\end{lemma}

\begin{proof}
The classical normally convergent expansion \cite[25.12.12]{dlmf-polylog}
gives, after pairing,
\begin{equation}\label{eq:deficit-pair}
 D_s(t):=\zeta(s)-\Li_s(e^{-t})
 =At-\frac{t^2}{2}K_\eps(-\Log t)+t^3 J_\eps(t),
\end{equation}
where $J_\eps$ is jointly holomorphic for $t$ near zero and $\eps$
near the indicated compact interval. Normal convergence can also be
checked by the zeta functional equation: the coefficients of the
analytic tail have geometric majorants on each smaller disk
$|t|<r<2\pi$. All poles near $\eps=0$ occur in the displayed paired term.

Using
\[
 \frac{1-e^{-\eps\ell}}{\eps}
  =\ell\int_0^1e^{-\eps\ell v}\,dv
\]
and the bounded regular part of $(E_2(\eps)-1)/\eps$ gives
\[
 |K_\eps(-\Log t)|
 \le C|t|^{-1/4}(1+|\log|t||)\le C'|t|^{-1/2}.
\]
Thus $D_s(t)=At+O(|t|^{3/2})=O(|t|)$. Uniformly in $b$,
\[
 b-\sqrt{b^2+2D}
 =-\frac Db+\frac{D^2}{2b^3}+O(|D|^3).
\]
Inserting \eqref{eq:deficit-pair} leaves the displayed terms of
\eqref{eq:local-exponent}; the difference between $D^2$ and
$A^2t^2$ is $O(|t|^{5/2})$. This proves the bound, including the
imaginary boundary rays by continuation and continuity.
\end{proof}

The nonlinear feedback contributes precisely $\nu$ to the quadratic
regular part. At criticality $b=A$, this is $1$, changing the
resonant constant from $3/2$ to $5/2$.

\subsection{A scalar equation absorbing the slow logarithms}

Define $B=B(n,\eps,b)>1$ by
\begin{equation}\label{eq:scale}
 \boxed{\displaystyle\quad
 V:=\frac{bB^2}{n}=K_\eps(\log B)+\nu.
 \quad}
\end{equation}
This defines a scale from the paired quadratic exponent; it does not
assume that the underlying distribution has a finite variance. Put
\begin{equation}\label{eq:q-xi}
 q=\frac{E_2(\eps)B^{-\eps}}{V},\qquad
 \xi_m=\frac{m-nA/b}{B},\qquad \xi=\xi_n.
\end{equation}

\begin{theorem}[Existence, uniqueness, and size of the scale]
\label{thm:scale}
Uniformly for $|\lambda|\le\Lambda$ and $b\in[b_-,b_+]$,
\eqref{eq:scale} has exactly one solution $B>1$ for all sufficiently
large $n$. Let
\begin{equation}\label{eq:leading-scale}
 \ph(\lambda)=\frac{1-e^{-\lambda}}{\lambda},\quad \ph(0)=1,
 \qquad B_0^2=\frac{nL}{2b}\ph(\lambda).
\end{equation}
Then
\begin{align}
 \frac B{B_0}
 &=1+\frac{C_\lambda}{L\ph(\lambda)}
       +O\left(\frac{(\log L)^2}{L^2}\right),\label{eq:scale-first}\\
 C_\lambda
 &=\gamma+\nu+e^{-\lambda}
   \left\{\frac12\log L+\frac12\log\frac{\ph(\lambda)}{2b}
                   -\gamma+\frac32\right\},\label{eq:C-lambda}\\
 q&=\frac{2e^{-\lambda}}{L\ph(\lambda)}
          +O\left(\frac{\log L}{L^2}\right).
          \label{eq:q-equivalent}
\end{align}
In particular $q>0$ and $q\asymp L^{-1}$ on the specified parameter sets.
\end{theorem}

\begin{proof}
Set $H=\log B$ and
$F(H)=be^{2H}/n-K_\eps(H)-\nu$. For small real $\eps$,
$K_\eps(0)+\nu$ is uniformly positive by
\eqref{eq:resonant-values}. Thus $F(0)<0$ for large $n$.
Moreover
\[
 F'(H)=e^{-\eps H}
 \left\{\frac{2b}{n}e^{(2+\eps)H}-E_2(\eps)\right\}.
\]
The braced expression increases strictly, because $2+\eps>0$.
It follows that $F$ first decreases and then increases, or increases
throughout. In both cases $F(H)\to+\infty$ as $H\to\infty$:
the growth exponent $2$ dominates the possible growth exponent
$-\eps$ of $K_\eps$. Hence there is exactly one positive zero.

Write $H_0=\log B_0$
$=\tfrac12\{L+\log L+\log(\ph(\lambda)/(2b))\}$.
The representation \eqref{eq:K}, with
$\eps=2\lambda/L$, gives uniformly, including $\lambda=0$,
\[
 K_\eps(H_0)+\nu
 =\frac L2\ph(\lambda)+C_\lambda
                   +O((\log L)^2/L).
\]
For example the first term follows from
$(1-e^{-\lambda})/\eps=L\ph(\lambda)/2$; the next term uses
$E_2'(0)=\gamma-3/2$. First bracket the root by $B=cB_0$
for fixed constants $c<1$ and $c>1$; the leading terms have opposite
signs. Then the derivative at the bracketed root is of order $L$.
Taylor's theorem in $H-H_0$ improves the bracket to
$H-H_0=C_\lambda/(L\ph(\lambda))+O((\log L)^2/L^2)$.
Exponentiation gives \eqref{eq:scale-first}. Finally
$E_2(\eps)B^{-\eps}=e^{-\lambda}(1+O(\log L/L))$
and $V=L\ph(\lambda)/2+O(\log L)$, proving the $q$ formula.
\end{proof}
